% References Figure~\ref{fig:highlights:lsstcam}, defined in the summary figure near the top of researchStatement.tex.
\textbf{Astronomical Instrumentation Development:}
My instrumentation work has encompassed both photometric and spectroscopic instruments, with a specific focus on focal plane technology.
This began with \LLAMAS, an integral field spectrograph for the Magellan Baade telescope. I validated the optical performance of the diffraction gratings and collimators for the \LLAMAS\ spectrographs, assembled and verified optical tolerance for the 24 spectrograph cameras, and integrated all 2,400 AR-coated fibers into the focal plane with zero fiber loss, before it was commissioned in 2024 \cite{Simcoe2022LLAMAS,llamas}. After integrating and testing \LLAMAS, I shifted my focus from integral-field spectrographs to multi-object spectrographs. While at the University of Michigan-Ann Arbor, I characterized and resolved a common failure mode on \DESI\ robotic positioners, known as the linear phi failure mode. While testing the \DESI\ positioners, I found that reorienting the robotic positioner to an inverted physical configuration could remove the linear phi behavior, indicating that gravity-sag has a significant influence on the failure mode \citep{Wuthrich2026Positioners}. % LLAMAS and DESI testing
% Recently, I have been testing robotic fiber positioners for next-generation spectroscopic focal planes, including the \must\ and the proposed \specsfive, and am currently leading fiber tilt-sensitivity tests of prototype positioners, including both static and dynamic measurements on a kinematic telescope mount. % Robotic positioners
This expertise with spectroscopic instruments translated effectively to the transcendent photometric imager of the next decade, the \LSSTCam. Since 2024, I have served as an \LSSTCam\ Summit Support Scientist, leading \LSSTCam\ focal-plane commissioning on Cerro Pach\'on as a key component of the in-kind contribution from Switzerland to the Rubin Project. I tested the CCDs in both the summit clean-room and on the Simonyi Survey Telescope, optimizing the sensor performance, and played a leading role in achieving \LSSTCam\ first photon in April 2024 (see Figure~\ref{fig:highlights:lsstcam}). The work to commission \LSSTCam, among other CCD optimizations, is detailed in the invited proceedings of SPIE Astronomical Telescopes + Instrumentation 2026, the most impactful astronomical instrumentation conference, held biennially \citep{MacBride2026LSSTCam}. % LSSTCam testing
% Additionally, I characterized CCD defects in the \LSSTCam\ focal plane, and developed a custom image-processing tool to monitor stability of defects during \LSSTCam operation. % LSSTCam
